\chapter{Kịch bản demo và minh chứng vận hành}
\label{app:demo}

Phụ lục này mô tả quy trình \emph{demo có thể tái lập} dùng để minh chứng hệ thống chạy
thật trên cụm: từ việc xác nhận eBPF can thiệp ở tầng nhân, mô phỏng tấn công, tới vòng
phát hiện~$\rightarrow$~phản hồi~$\rightarrow$~cách ly~$\rightarrow$~xác minh, cùng cách
truy cập kho đối tượng MinIO và cách làm cho dashboard quan trắc hiển thị dữ liệu. Mọi
lệnh đều chạy trên cụm thật ($3$~node: \texttt{userver-master}, \texttt{userver-worker-1},
\texttt{userver-worker-home}).

\section{Chuẩn bị và xác nhận hiện trạng cụm}

\begin{lstlisting}[language={},caption={Xác nhận node và các pod thành phần đang chạy.}]
$ kubectl get nodes
$ kubectl -n zt-mapper get pods -o wide   # DaemonSet thu thap/pipeline tren moi node
$ kubectl -n soc       get pods           # NATS, postgres, aggregator, incident, console, minio
\end{lstlisting}

DaemonSet \texttt{zt-pipeline-ebpf} hiện diện trên cả ba node (mỗi pod gồm ba container
\texttt{collector}/\texttt{pipeline}/\texttt{alert-bridge}); các dịch vụ SOC nằm trong
không gian tên \texttt{soc}.

\section{Minh chứng eBPF can thiệp ở tầng nhân}
\label{app:demo-ebpf}

Đây là phần cốt lõi cần thể hiện khi bảo vệ. Mục tiêu: chứng minh chương trình eBPF đã
\emph{nạp vào nhân}, \emph{qua trình kiểm chứng}, được \emph{JIT}, \emph{gắn móc} vào
\texttt{tcp\_v4\_connect}, và đang \emph{bắt sự kiện} theo thời gian thực. Vì
\texttt{DaemonSet} thu thập chạy \texttt{hostNetwork} và đặc quyền, chương trình eBPF nằm
trong nhân của \emph{host}; do đó \texttt{bpftool} được chạy \emph{trên chính node} (đăng
nhập node qua SSH, hoặc dùng container gỡ lỗi mức node):

\begin{lstlisting}[language={},caption={Bốn minh chứng eBPF ở tầng nhân (chạy trên node, chụp màn hình từng bước).}]
# Dang nhap node worker-1 (noi dat pod tan cong) roi chay:

# (1) Chuong trinh da nap: kieu kprobe, gan vao tcp_v4_connect, da qua verifier + JIT
$ sudo bpftool prog show | grep -A4 -i 'kprobe\|tcp_v4_connect'

# (2) Ba BPF map: cidr_map (lpm_trie), events (ringbuf), sock_store (hash)
$ sudo bpftool map show | grep -i 'cidr_map\|events\|sock_store'

# (3) Doi chieu: cac CIDR pod-network da nap vao bo loc LPM
$ sudo bpftool map dump name cidr_map | head

# (4) Tracing thoi gian thuc tu trong container (khong can bpftool):
#     moi connect() sinh mot su kien JSON
$ POD=$(kubectl -n zt-mapper get pod -l app.kubernetes.io/name=zt-pipeline-ebpf \
        --field-selector spec.nodeName=userver-worker-1 \
        -o jsonpath='{.items[0].metadata.name}')
$ kubectl -n zt-mapper logs $POD -c collector -f
\end{lstlisting}

\noindent\textit{[Chèn ảnh chụp đầu ra của bốn lệnh trên~--~đây là minh chứng mạnh nhất
cho phần eBPF.]} Điểm cần chỉ ra khi trình bày: dòng \texttt{kprobe ... tcp\_v4\_connect}
(điểm móc), các trường \texttt{xlated/jited} (đã qua verifier và JIT), và ba map đúng
kiểu \texttt{lpm\_trie}/\texttt{ringbuf}/\texttt{hash}. So với một agent ở không gian
người dùng, đây là bằng chứng hệ thống quan sát \emph{tại nguồn} trong nhân với chi phí
cố định trên mỗi lời gọi.

\section{Mô phỏng tấn công: quét cổng và lan truyền ngang}
\label{app:demo-attack}

Một pod đóng vai kẻ tấn công trong không gian tên \texttt{zt-targets} thực hiện quét cổng
\texttt{connect} (\texttt{nmap -sT}) tới nhiều pod dịch vụ, vượt ngưỡng $\tau=40$ kết nối
trong cửa sổ thời gian để kích hoạt cảnh báo, đồng thời chạm tới nhiều đích để kích hoạt
cảnh báo \emph{blast radius}.

\begin{lstlisting}[language={},caption={Kịch bản tấn công từ pod \texttt{attacker}.}]
# Quet cong connect-scan toi dai pod dich vu (vuot nguong -> port_scan)
$ kubectl -n zt-targets exec attacker -- \
    nmap -sT -p 1-200 --max-rate 200 10.244.204.0/24

# Cham toi nhieu dich de kich hoat blast_radius (lan truyen ngang)
$ kubectl -n zt-targets exec attacker -- \
    sh -c 'for ip in 10.244.204.{150..180}; do nc -w1 -z $ip 80; done'
\end{lstlisting}

\section{Vòng phát hiện -- phản hồi -- cách ly -- xác minh}

\begin{lstlisting}[language={},caption={Quan sát vòng end-to-end từ phát hiện tới cách ly.}]
# (1) Pipeline phat canh bao (port_scan / blast_radius)
$ kubectl -n zt-mapper logs $POD -c pipeline | grep ALERT

# (2) Dich vu su co tao su co + tu sinh ban nhap NetworkPolicy
$ kubectl -n soc logs deploy/zt-incident-service | grep -E 'incident.created|action'

# (3) Phe duyet tren console (console.zt.local) -> hanh dong chuyen 'executed'
#     -> NetworkPolicy cach ly duoc ap
$ kubectl -n zt-targets get networkpolicy | grep quarantine

# (4) Xac minh hieu luc: pod tan cong mat ket noi di ra
$ kubectl -n zt-targets exec attacker -- nc -w2 10.244.204.161 80; echo "exit=$?"
# exit=1  (truoc khi cach ly: exit=0)
\end{lstlisting}

\noindent\textit{[Chèn ảnh chụp console: hàng đợi phê duyệt, chi tiết sự cố, và trạng
thái \texttt{executed} của hành động cách ly~--~xem các Hình ở
mục~\ref{sec:impl-app}.]}

\section{Truy cập kho đối tượng MinIO}
\label{app:demo-minio}

MinIO (không gian tên \texttt{soc}) là kho đối tượng lưu báo cáo tuần, bản sao lưu cơ sở
dữ liệu và bản nháp \texttt{NetworkPolicy}. MinIO \emph{không} mở ra ngoài qua ingress;
truy cập an toàn bằng cách chuyển tiếp cổng tới giao diện console (\texttt{:9001}):

\begin{lstlisting}[language={},caption={Truy cập MinIO console và liệt kê bucket.}]
# Chuyen tiep cong console MinIO ra may cuc bo
$ kubectl -n soc port-forward svc/minio 9001:9001
# Mo http://localhost:9001  (user: zt-soc-admin)

# Hoac dung mc qua port-forward API (:9000)
$ kubectl -n soc port-forward svc/minio 9000:9000 &
$ mc alias set zt http://localhost:9000 zt-soc-admin '<minio-password>'
$ mc ls zt   # zt-postgres-backups, zt-reports, zt-pcap-dumps, zt-netpol-archives
\end{lstlisting}

Trên trang \emph{Reports} của console, nút \emph{Generate \& upload to MinIO} kết xuất
báo cáo tuần và đẩy lên bucket \texttt{zt-reports}; có thể kiểm chứng bằng
\texttt{mc ls zt/zt-reports}.

\section{Làm cho dashboard quan trắc hiển thị dữ liệu}
\label{app:demo-grafana}

Console/Reports đọc trực tiếp PostgreSQL nên luôn hiển thị đầy đủ sự cố tích lũy. Trái
lại, dashboard Grafana là góc nhìn theo cửa sổ thời gian và phụ thuộc nguồn thu thập, nên
có thể trống dù cơ sở dữ liệu đầy. Để dashboard hiển thị dữ liệu cần:

\begin{enumerate}
  \item \textbf{Bảng \emph{Alerts Overview} (nguồn Loki).} Nới cửa sổ thời gian (mặc định
        $30$~phút) lên phạm vi bao trùm đợt tấn công, hoặc chạy lại kịch bản tấn công
        ở mục~\ref{app:demo-attack} ngay trước khi xem để có nhật ký cảnh báo trong cửa sổ.
  \item \textbf{Bảng \emph{Incidents \& Actions} (nguồn Prometheus).} Các bộ đếm
        (\texttt{zt\_incidents\_created\_total}, \texttt{zt\_aggregator\_active\_incidents}
        \ldots) chỉ hiển thị khi Prometheus thật sự thu thập điểm \texttt{/metrics} của
        dịch vụ SOC. Cụm dùng Prometheus Operator (kube-prometheus-stack), vốn khám phá
        mục tiêu qua \texttt{ServiceMonitor} chứ không qua chú thích
        \texttt{prometheus.io/scrape}; vì vậy cần khai báo một \texttt{ServiceMonitor} cho
        các dịch vụ SOC (kèm nhãn \texttt{release} mà Operator chọn). Ngoài ra các bộ đếm
        là biến đếm trong tiến trình~--~bị đặt lại $0$ mỗi khi pod khởi động lại~--~nên
        chúng phản ánh hoạt động \emph{kể từ lần khởi động gần nhất}, không phải lịch sử.
\end{enumerate}

\noindent Đây cũng là một giới hạn quan trắc được bàn ở Chương~\ref{ch:eval}: cơ sở dữ
liệu là nguồn dữ liệu bền vững và là hệ thống ghi nhận chính thức của các số liệu sự cố;
dashboard là lớp quan trắc thời gian thực bổ trợ.
